\clearpage
\VerdierPhysicalPageAnchor{0229}%
\begingroup\small\setlength{\parskip}{0pt}\setlength{\emergencystretch}{1.5em}
\setlength{\abovedisplayskip}{4pt plus 2pt minus 2pt}%
\setlength{\belowdisplayskip}{4pt plus 2pt minus 2pt}%

\VerdierSemanticAnchor{ch4d-4-1-6-cont}
\noindent the $n$-fold complex defined by:

\VerdierSemanticAnchor{ch4d-4-1-6-1}
\begin{equation}
({}^{m}\!X(p,q))^\sigma=
\begin{cases}
X^\sigma,&-q\leq\sigma_m<-p,\\
0,&\sigma_m<-q\ \lor\ -p\leq\sigma_m,
\end{cases},
\tag{4.1.6.1}
\end{equation}

\noindent the differentials of ${}^{m}\!X(p,q)$ are the differentials of $X$,
except possibly when the source or target of the differential under
consideration is a zero object.

For every pair $(p,q)$ and $(p',q')$ of elements of
$\mathcal{F}\ell_1(\widetilde{\mathbb Z})$ such that $p\leq p'$ and
$q\leq q'$, denote by

\VerdierSemanticAnchor{ch4d-4-1-6-2}
\begin{equation}
f\bigl((p,q),(p',q')\bigr)\colon
{}^{m}\!X(p,q)\longrightarrow{}^{m}\!X(p',q')
\tag{4.1.6.2}
\end{equation}

the morphism of $n$-fold complexes whose components are identity
morphisms, except possibly when the source or target is a zero object.
Finally, for every element $(p,q,r)$ of
$\mathcal{F}\ell_2(\widetilde{\mathbb Z})$, denote by

\VerdierSemanticAnchor{ch4d-4-1-6-3}
\begin{equation}
g(p,q,r)\colon{}^{m}\!X(q,r)
\longrightarrow T_m\bigl({}^{m}\!X(p,q)\bigr)
\tag{4.1.6.3}
\end{equation}

($T_m$ denotes the translation functor in direction $m$
(Chapter~I, 2.3.1)) the morphism whose components are:

\VerdierSemanticAnchor{ch4d-4-1-6-4}
\begin{equation}
\left\{
\begin{aligned}
g(p,q,r)^\sigma&=0,
  &&\sigma_m\neq-q-1,\\
g(p,q,r)^\sigma&=-d_X^{\sigma,m},
  &&\sigma_m=-q-1.
\end{aligned}
\right.
\tag{4.1.6.4}
\end{equation}

We now introduce new notation. First, by an abuse of notation, denote by
$\int_n\colon\operatorname{comp}^n(\mathcal A)\longrightarrow
\operatorname{comp}(\mathcal A)$ the $\Pi$-extension or the
$\Sigma$-extension of the associated simple-complex functor (provided
that countable products or sums exist in $\mathcal A$
(Chapter~I, 2.2.3)). Set:

\VerdierSemanticAnchor{ch4d-4-1-6-5}
\begin{equation}
\left\{
\begin{aligned}
\int_n{}^{m}\!X(p,q)
  &={}^{m}\!Y(p,q),\\
\int_n{}^{m}\!X(-\infty,q)
  &={}^{m}\!Y(-\infty,q)={}^{m}\!Y(q),\\
\int_n{}^{m}\!X(-\infty,+\infty)=\int_n X
  &={}^{m}\!Y(-\infty,+\infty)={}^{m}\!Y(+\infty)=Y.
\end{aligned}
\right.
\tag{4.1.6.5}
\end{equation}
\endgroup
